\documentclass[11pt]{article}

% ---------- Packages ----------
\usepackage[margin=1in]{geometry}
\usepackage{amsmath,amssymb,amsthm,mathtools}
\usepackage{bm}
\usepackage{graphicx}
\usepackage{booktabs}
\usepackage[hidelinks]{hyperref}
\usepackage{microtype}
\usepackage{enumitem}
\usepackage{caption}
\usepackage{subcaption}
\usepackage[ruled,vlined]{algorithm2e}

% --- If you prefer biblatex, uncomment below and comment out natbib in your toolchain ---
% \usepackage[backend=biber,style=authoryear,maxcitenames=2]{biblatex}
% \addbibresource{refs.bib}

% ---------- Macros ----------
\newcommand{\R}{\mathbb{R}}
\newcommand{\E}{\mathbb{E}}
\newcommand{\KL}{\mathrm{D}_{\mathrm{KL}}}
\newcommand{\Pbb}{\mathbb{P}}
\newcommand{\Var}{\mathrm{Var}}
\newcommand{\tr}{\mathrm{tr}}
\newcommand{\T}{\top}
\newcommand{\Khat}{\widehat{K}}
\newcommand{\nuY}{\nu}
\newcommand{\Gammahat}{\widehat{\Gamma}}
\newcommand{\Y}{\mathcal{Y}}

% ---------- Title ----------
\title{Least-Information Action on Delay Embeddings: \\
Testing Rose--Takens with Local Linear Diffusion Kernels}
\author{%
  % TODO: Authors
}
\date{\today}

\begin{document}
\maketitle

\begin{abstract}
\noindent
% TODO (120--150 words): one-sentence problem; least-information (edge-KL) action; local linear + residual diffusion kernel; Rose test $\Rightarrow$ empirical evidence for measure + spectral equivalence; brief results line.
\end{abstract}

% =========================================================
\section{Introduction}
% One-paragraph gap: Takens (geometry) vs statistics/spectra.
% One-paragraph idea: minimize one-step KL "action" between empirical and model edge laws.
% Bulleted contributions: kernel estimator; KL-action test; optional spectral/geometry checks; basis selection note.

\paragraph{Thesis (applied).}
Choose coordinates and a local diffusion kernel that minimize the one-step KL \emph{edge action} between data and model; when the action is $\approx 0$ (the Rose condition), the reconstruction attains measure and spectral equivalence.

% =========================================================
\section{Setup and Principle}
\label{sec:setup}
Let $x_t$ evolve on an attractor $A$; form a delay vector $y_t=\Psi_k(x_t)\in\R^{k+1}$.
Optionally reduce with a linear basis $z_t=Ly_t\in\R^{m}$, $m\ge 2d+1$.
Let $\nu$ be the stationary law on $\Y$ (either $Y$ or $Z$ space).

\paragraph{Kernel and edge law.}
A one-step transition kernel $K(y_0,\mathrm{d}y_1)$ induces the joint edge measure
\begin{equation}
\Gamma(\mathrm{d}y_0,\mathrm{d}y_1)=\nu(\mathrm{d}y_0)\,K(y_0,\mathrm{d}y_1).
\end{equation}

\paragraph{Least-information action.}
Given a fitted model kernel $\Khat$, define the \emph{edge action}
\begin{equation}
\label{eq:edge-action}
\mathcal{A}_{\text{edge}}
\;=\;
\KL\!\big(\Gamma_{\text{data}} \,\|\, \Gamma_{\text{model}}\big)
\;=\;
\KL\!\big(\nu(\mathrm{d}y_0)K\,\|\,\nu(\mathrm{d}y_0)\Khat\big).
\end{equation}
The \emph{Rose test} accepts when $\mathcal{A}_{\text{edge}}\approx 0$ (with uncertainty), evidencing kernel equality on support.

% =========================================================
\section{Methods}
\label{sec:methods}

\subsection{Basis (optional) and neighborhoods}
\label{subsec:basis}
% Keep it minimal; you can expand later.
We over-embed via $\Psi_k$ and optionally choose a well-conditioned linear basis $L$ (e.g., PCA-whiten $Y_t$ then apply a small random rotation; or use a Fisher/CCA basis). Neighborhoods around a query state $y_0$ (or $z_0$) are built with kNN or kernel weights $w_t=w(\|y_t-y_0\|/h)$.

\subsection{Local linear diffusion kernel}
\label{subsec:kernel}
For a query $y_0$, fit a locally weighted linear map
\begin{equation}
\label{eq:local-mean}
(A(y_0),b(y_0)) \;=\; \arg\min_{A,b} \sum_t w_t\,\|y_{t+1}-A y_t - b\|^2,
\qquad
\mu(y_0)=A(y_0)\,y_0+b(y_0).
\end{equation}
Residual covariance (ridge-stabilized):
\begin{equation}
\label{eq:local-cov}
\Sigma(y_0)
=
\frac{\sum_t w_t\, r_t r_t^\T}{\sum_t w_t}
+\lambda I,
\quad
r_t=y_{t+1}-A(y_0) y_t - b(y_0).
\end{equation}
Define the empirical kernel (single-Gaussian; mixture optional):
\begin{equation}
\label{eq:kernel}
\Khat(y_0,\mathrm{d}y_1)=\mathcal{N}\!\big(\mu(y_0),\Sigma(y_0)\big)\,\mathrm{d}y_1.
\end{equation}

\subsection{Hyperparameters}
\label{subsec:hyper}
Choose $h$ (neighborhood scale), $\lambda$ (ridge), and optional mixture size by maximizing held-out conditional log-likelihood
$\sum \log \Khat(y_t,y_{t+1})$ (equivalently, minimizing forward KL).

\subsection{Estimating the edge action (Rose test)}
\label{subsec:rose}
Use a time-blocked hold-out set $\mathcal{H}$ of edges to estimate
\begin{equation}
\widehat{\mathcal{A}}_{\text{edge}}
=
\frac{1}{|\mathcal{H}|}
\sum_{(y_0,y_1)\in\mathcal{H}}
\log\frac{\hat p_{\text{emp}}(y_1\mid y_0)}{\Khat(y_0,y_1)},
\end{equation}
where $\hat p_{\text{emp}}(\cdot\mid y_0)$ is a light-smoothing baseline (e.g., coarse bins or kNN) ensuring domination. Report a bootstrap CI; accept Rose if CI contains $0$. Optionally report a symmetric proxy (JS divergence).

\subsection{Optional: Spectral/geometry diagnostics}
\label{subsec:spectral}
(i) Koopman correlation moments $c_k=\langle h, U^k h\rangle$ vs empirical; 
(ii) Diffusion maps from a symmetrized affinity
$W(y_i,y_j)=\tfrac12\big[\Khat(y_i,y_j)\nu(y_i)+\Khat(y_j,y_i)\nu(y_j)\big]$
to compare leading spectra.

% =========================================================
\section{Algorithms (concise)}
\label{sec:algs}

\begin{algorithm}[H]
\DontPrintSemicolon
\SetAlgoLined
\caption{FitLocalDiffusionKernel($y_0$, data, $h$, $\lambda$)}
Compute weights $w_t=w(\|y_t-y_0\|/h)$.\;
Solve \eqref{eq:local-mean} for $(A,b)$; set $\mu(y_0)=A y_0+b$.\;
Compute $\Sigma(y_0)$ via \eqref{eq:local-cov}.\;
Return $\Khat(y_0,\cdot)=\mathcal{N}(\mu(y_0),\Sigma(y_0))$.\;
\end{algorithm}

\begin{algorithm}[H]
\DontPrintSemicolon
\SetAlgoLined
\caption{RoseTest(data, $\Khat$)}
Hold out a time block $\mathcal{H}$; build $\hat p_{\text{emp}}(\cdot\mid y_0)$.\;
Compute $\widehat{\mathcal{A}}_{\text{edge}}$ and a bootstrap CI.\;
\textbf{Decision:} accept Rose if CI contains $0$.\;
\end{algorithm}

% =========================================================
\section{Experiments}
\label{sec:expts}
% Brief, task-focused template; expand with specifics later.

\paragraph{Datasets.}
% TODO: logistic, Hénon, Lorenz (discrete sample); optional real series.

\paragraph{Protocol.}
Over-embed $\Psi_k$; optional basis $L$; time-blocked train/val/test; prevent leakage.

\paragraph{Baselines.}
Ulam counts (partition), kNN/KDE conditional.

\paragraph{Metrics.}
KL-action (and CI), predictive NLL, optional spectral/DMAPS checks.

% =========================================================
\section{Results}
\label{sec:results}
% Tables/figures with minimal commentary. Expand later.
% TODO: Figure placeholders
% \begin{figure}[h!]
%   \centering
%   \includegraphics[width=.48\textwidth]{fig_dmaps.pdf}
%   \includegraphics[width=.48\textwidth]{fig_kl_vs_h_lambda.pdf}
%   \caption{Left: DMAPS; Right: KL-action vs hyperparameters.}
% \end{figure}

% =========================================================
\section{Discussion}
% When local Gaussian suffices; interpreting Rose failure (insufficient basis, sample size, misspecification); deterministic (vanishing-viscosity) note.

% =========================================================
\section{Conclusion}
% One-paragraph recap: least-information action; Rose test; evidence for measure+spectral equivalence; quick guidance.

% =========================================================
\paragraph{Reproducibility.}
% Code/data availability statement; seeds; config table (appendix).

% =========================================================
% \printbibliography % (if using biblatex)

\appendix

\section{Least-Information Action (reference)}
\label{app:action}
Path-space variant (optional): for a Markov model $Q$ vs data/reference $P$,
\begin{equation}
\KL(\Pbb_Q\|\Pbb_P)
=
\KL(Q_0\|P_0)
+\sum_{t=0}^{T-1}\E_Q\!\Big[\log \tfrac{K_Q(Y_t,Y_{t+1})}{K_P(Y_t,Y_{t+1})}\Big].
\end{equation}
Near a minimum, the Fisher--Rao quadratic supplies a local metric:
$\KL(\theta+\delta\theta\|\theta)=\tfrac12\,\delta\theta^\T I(\theta)\,\delta\theta+o(\|\delta\theta\|^2)$.

\end{document}
